
本附录将自足地证明：带\BookRecordIndex{T0066}标记点的亏格零曲线\BookRecordIndex{T0246}模空间的 Grothendieck--Knudsen 紧化是复代数流形。

D.1 节首先预备性地讨论黎曼球面上的\BookRecordIndex{T0185}交比与\BookRecordIndex{T0126}分式线性变换。D.2 节研究顶点由
$1$ 到 $n$ 的整数标号的稳定树 $T$。每个整数恰出现一次；一个顶点可以带多个标号，
也可以不带标号，但每个顶点处标号数与边数之和至少为三。切断一条边使指标集
$\{1,\ldots,n\}$ 分割为两个子集；我们将说明如何从相关的\BookRecordIndex{T0410}指标集分割集合 $S$
恢复树的同构类。

D.3 节主要由定义组成，引入带 $n$ 个\BookRecordIndex{T0066}标记点的亏格零稳定黎曼曲线的\BookRecordIndex{T0246}模空间
$\overline{\mathcal M}_{0,n}$。D.4 节说明如何利用由\BookRecordIndex{T0185}交比定义的函数
$w_{ijk\ell}:\overline{\mathcal M}_{0,n}\to S^2$，将
$\overline{\mathcal M}_{0,n}$ 嵌入二球面的乘积。D.4 节还证明此映射的像
$\overline{\mathcal M}_{0,n}$ 是 $(S^2)^N$ 的光滑子流形，且每棵\BookRecordIndex{T0084}带标号树确定一个
光滑子流形（Knudsen）。Keel 证明了这些子流形生成
$\overline{\mathcal M}_{0,n}$ 的上同调，并计算了环结构。D.5 节考察
$\overline{\mathcal M}_{0,n}$ 的拓扑，D.6 节讨论其泛性质，D.7 节说明它与
\BookRecordIndex{T0039}Mumford 商的关系。

我们通过\BookRecordIndex{T0185}交比构造 Grothendieck--Knudsen 紧化的方法显然是新的。所用数学技巧虽初等
却精妙。把 $\overline{\mathcal M}_{0,n}$ 嵌入二球面乘积，让人联想到 Fulton 与
MacPherson 对更高维簇的构造（参见~\srcite{132}）。扩展\BookRecordIndex{T0185}交比
$w_{ijk\ell}:\overline{\mathcal M}_{0,n}\to S^2$ 也出现在 Keel~\srcite{208} 的工作中。

\section{\BookRecordIndex{T0040}Möbius 变换与\BookRecordIndex{T0185}交比}\label{sec:D-1}

全文通过从北极出发的球极投影
\[
  S^2\longrightarrow\mathbb C\cup\{\infty\}:\quad
  x\longmapsto z=\frac{x_1+ix_2}{1-x_3}
\]
把单位球面 $S^2\subset\mathbb R^3$ 与黎曼球面 $\mathbb C\cup\{\infty\}$ 等同。
此微分同胚把两种标准复结构（$S^2$ 上的 $\xi\mapsto x\times\xi$ 与
$\mathbb C$ 上的 $i$）等同，并把 $S^2$ 上标准度量的四分之一与
$\mathbb C\cup\{\infty\}$ 上的 Fubini--Study 度量等同（见习题~\ref{ex:4-2-1}）。
相应体积形式为
\[
  \mathrm{dvol}=\frac{dx\wedge dy}{(1+x^2+y^2)^2},
\]
其中 $z=x+iy$。在这一度量下，二球面的面积为 $\pi$。

\noindent\textbf{\BookRecordIndex{T0040}Möbius 变换。}
\BookRecordIndex{T0040}Möbius 变换是二球面的全纯微分同胚。这种微分同胚具有形式
\[
  \phi(z)=\frac{az+b}{cz+d},\qquad ad-bc=1,
\]
也称为\emph{\BookRecordIndex{T0126}分式线性变换}。全文以 $G=\operatorname{PSL}(2,\mathbb C)$ 表示
\BookRecordIndex{T0040}Möbius 变换群。等距变换由条件 $d=\bar a$、$c=-\bar b$ 刻画，构成子群
\[
  \operatorname{SO}(3)=\operatorname{SU}(2)/\{\pm 1\}\subset G.
\]
这是 $G$ 的一个极大紧子群。

\noindent\textbf{\BookRecordIndex{T0185}交比。}
给定 $S^2=\mathbb C\cup\{\infty\}$ 中任意三个互异点 $z_0,z_1,z_2$，
存在唯一的 \BookRecordIndex{T0040}Möbius 变换 $\phi\in G$ 将它们依次映到 $0,1,\infty$。它由
$z\mapsto w(z_0,z_1,z_2,z)$ 给出，其中
\[
  w(z_0,z_1,z_2,z_3):=
  \frac{(z_1-z_2)(z_3-z_0)}{(z_0-z_1)(z_2-z_3)}.
\]
这个数称为 $z_0,z_1,z_2,z_3$ 的\emph{\BookRecordIndex{T0185}交比}。它对 $S^2$ 中每个四点互异的
四元组都有定义，且在该定义域上取值于 $S^2\setminus\{0,1,\infty\}$。
注意 $w(0,1,\infty,z)=z$。此外，$w$ 可连续延拓到所有至多两点相等的
四元组组成的集合；不允许三个点重合。因此
$w:(S^2)^4\setminus\Delta_3\to S^2$，其中
\[
  \Delta_3:=\{(z_0,z_1,z_2,z_3)\mid
    \exists\, i<j<k\text{ 使 }z_i=z_j=z_k\}.
\]
这一延拓具有以下性质。

\noindent\textbf{（分式线性性）}
在 $z_0,z_1,z_2,z_3$ 中任取三个互异点，\BookRecordIndex{T0185}交比作为第四个点的函数是\BookRecordIndex{T0126}分式线性变换。

\noindent\textbf{（不变性）}
对 $\phi\in G$ 与 $(z_0,z_1,z_2,z_3)\in(S^2)^4\setminus\Delta_3$，
\[
  w(\phi(z_0),\phi(z_1),\phi(z_2),\phi(z_3))=w(z_0,z_1,z_2,z_3).
\]

\noindent\textbf{（对称性）}
对 $\{i,j,k,\ell\}=\{0,1,2,3\}$ 的任意排列，简记
$w_{ijk\ell}:=w(z_i,z_j,z_k,z_\ell)$。则
\begin{equation}\label{eq:D-1-1}
  w_{jik\ell}=w_{ij\ell k}=1-w_{ijk\ell},\qquad
  w_{ikj\ell}=\frac{w_{ijk\ell}}{w_{ijk\ell}-1}.
\end{equation}

\noindent\textbf{（归一化）}
对 $(z_0,z_1,z_2,z_3)\in(S^2)^4\setminus\Delta_3$，
\begin{equation}\label{eq:D-1-2}
  w_{0123}=\begin{cases}
    \infty,&z_0=z_1\text{ 或 }z_2=z_3,\\
    1,&z_0=z_2\text{ 或 }z_3=z_1,\\
    0,&z_0=z_3\text{ 或 }z_1=z_2.
  \end{cases}
\end{equation}

\noindent\textbf{（递推性）}
若 $(z_0,z_1,z_2,z_3)$、$(z_0,z_1,z_2,z_4)$ 和
$(1,\infty,w_{0123},w_{0124})$ 都属于 $(S^2)^4\setminus\Delta_3$，则
\begin{equation}\label{eq:D-1-3}
  w_{1234}=\frac{w_{0124}-1}{w_{0124}-w_{0123}}.
\end{equation}

\begin{exercise}\label{ex:D-1-1}
证明上述递推性。提示：若 $z_0,z_1,z_2$ 互异，则存在唯一的 \BookRecordIndex{T0040}Möbius 变换，
把 $z_0,z_1,z_2,z_3,z_4$ 映到 $0,1,\infty,u,v$，其中
$u:=w_{0123}$、$v:=w_{0124}$。若 $z_1,z_2,z_3$ 互异，
则 $1,\infty,u$ 也互异，且把 $(1,\infty,u)$ 映到 $(0,1,\infty)$ 的
唯一元素 $\phi\in G$ 具有形式 $\phi(z)=(z-1)/(z-u)$。
故 $w_{1234}=w(1,\infty,u,v)=\phi(v)=(v-1)/(v-u)$。
\end{exercise}

\noindent\textbf{\BookRecordIndex{T0243}冒泡。}
群 $G$ 非紧。由习题~\ref{ex:4-2-3}，导数一致有界的任意序列
$\phi^\nu\in G$ 都有子序列在整个 $S^2$ 上连同所有阶导数一致收敛。
注意序列 $\phi^\nu(z):=\nu z+\nu^2$ 不存在这样的子序列，尽管对每个
$z\in S^2$ 都有 $\phi^\nu(z)\to\infty$。在每个有界集合
$\{|z|\le c\}$ 上，收敛是一致的；但在 $z=\infty$ 的任何邻域上都不是
一致的，因为 $\phi^\nu(-\nu)=0$。下面将刻画 $G$ 中没有收敛子序列的序列。

\begin{lemma}\label{lem:D-1-2}
设 $\phi^\nu:S^2\to S^2$ 是一列 \BookRecordIndex{T0040}Möbius 变换，没有一致收敛的子序列。
则存在 $x,y\in S^2$ 和子序列 $\phi^{\nu_i}$，使其在
$S^2\setminus\{x\}$ 的紧子集上\BookRecordIndex{T0416}逐紧一致收敛于 $y$。
\end{lemma}
\begin{proof}
选择 $\rho^\nu\in\operatorname{SO}(3)$，使
$\rho^\nu\circ(\phi^\nu)^{-1}(\infty)=\infty$。
取子序列后，可假设 $\rho^\nu$ 在整个 $S^2$ 上一致收敛于 $\rho$。
序列 $\psi^\nu:=\phi^\nu\circ(\rho^\nu)^{-1}$ 满足
$\psi^\nu(\infty)=\infty$，因而形如
\[
  \psi^\nu(z)=a^\nu z+b^\nu,\qquad a^\nu\ne0.
\]
必要时把 $\psi^\nu$ 换成 $(\psi^\nu)^{-1}$，\sourcefootnote{原书此处允许改用逆序列，但末尾仍用 $\phi^\nu=\psi^\nu\circ\rho^\nu$ 回推结论；若已改用逆序列，该等式便不再成立。需另行说明逆序列的\BookRecordIndex{T0416}逐紧一致收敛如何转回原序列，参见习题 D.1.3。}可假设 $a^\nu$ 有界；
再取子序列，可假设 $a^\nu\to a\in\mathbb C$。进一步取子序列，
可假设或者 $b^\nu\to b\in\mathbb C$，或者 $|b^\nu|\to\infty$。
若 $|b^\nu|\to\infty$，则 $\psi^\nu$ 在
$\mathbb C=S^2\setminus\{\infty\}$ 的紧子集上\BookRecordIndex{T0416}逐紧一致收敛于 $\infty$。
若 $b^\nu\to b$，则 $a^\nu\to0$，从而 $\psi^\nu$ 在
$\mathbb C$ 的紧子集上\BookRecordIndex{T0416}逐紧一致收敛于 $b$。两种情形下，$\psi^\nu$
都满足引理的结论，因此 $\phi^\nu=\psi^\nu\circ\rho^\nu$ 也满足。
这证明了引理~\ref{lem:D-1-2}。
\end{proof}

若 $\phi^\nu$ 在 $S^2\setminus\{x\}$ 的紧子集上\BookRecordIndex{T0416}逐紧一致收敛于 $y$，
我们说序列 $\phi^\nu$ 在 $S^2\setminus\{x\}$ 上\emph{\BookRecordIndex{T0416}逐紧一致收敛}
于 $y$。注意，这不对序列 $\phi^\nu(x)$ 的行为作任何断言。

\begin{exercise}\label{ex:D-1-3}
设 $\phi^\nu\in G$ 在 $S^2\setminus\{x\}$ 上\BookRecordIndex{T0416}逐紧一致收敛于
$y\in S^2$。对每个 $w\in S^2$，给出使 $\phi^\nu(x)\to w$ 的例子。
还要证明，在所有情形下，$(\phi^\nu)^{-1}$ 在
$S^2\setminus\{y\}$ 上\BookRecordIndex{T0416}逐紧一致收敛于 $x$。
\end{exercise}

以下技术引理将在 D.5 节和第 5 章中使用。

\begin{lemma}\label{lem:D-1-4}
设 $\phi^\nu:S^2\to S^2$ 是一列 \BookRecordIndex{T0040}Möbius 变换。假设
$x_0,y_0\in S^2$，且 $x_1^\nu,x_2^\nu,y^\nu\in S^2$
为收敛序列，满足
\[
  x_0\ne\lim_{\nu\to\infty}x_1^\nu
  \ne\lim_{\nu\to\infty}x_2^\nu\ne x_0,
  \qquad y_0\ne\lim_{\nu\to\infty}y^\nu,
\]
以及
\[
  \lim_{\nu\to\infty}\phi^\nu(x_1^\nu)
  =\lim_{\nu\to\infty}\phi^\nu(x_2^\nu)=y_0,
  \qquad
  \lim_{\nu\to\infty}(\phi^\nu)^{-1}(y^\nu)=x_0.
\]
则 $\phi^\nu$ 在 $S^2\setminus\{x_0\}$ 的紧子集上\BookRecordIndex{T0416}逐紧一致收敛于 $y_0$。
\end{lemma}
\begin{proof}
记 $y:=\lim_{\nu\to\infty}y^\nu$，以及
$x_i:=\lim_{\nu\to\infty}x_i^\nu$（$i=1,2$）。选取等距变换序列
$\sigma^\nu,\rho^\nu\in\operatorname{SO}(3)\subset G$，它们都一致收敛于恒等变换，
并满足
\[
  \rho^\nu(y^\nu)=y,\qquad
  \sigma^\nu\circ(\phi^\nu)^{-1}(y^\nu)=x_0.
\]
记 $\xi_1^\nu:=\sigma^\nu(x_1^\nu)$、
$\xi_2^\nu:=\sigma^\nu(x_2^\nu)$。于是序列
\[
  \psi^\nu:=\rho^\nu\circ\phi^\nu\circ(\sigma^\nu)^{-1}
\]
满足
\[
  \psi^\nu(x_0)=y,\qquad
  \lim_{\nu\to\infty}\psi^\nu(\xi_1^\nu)
  =\lim_{\nu\to\infty}\psi^\nu(\xi_2^\nu)=y_0.
\]
不失一般性，可假设 $y_0=0$ 且 $x_0=y=\infty$。此时
$\psi^\nu(\infty)=\infty$，故对某些 $a^\nu,b^\nu\in\mathbb C$，
\[
  \psi^\nu(z)=a^\nu z+b^\nu.
\]
这些序列满足
\[
  \lim_{\nu\to\infty}(a^\nu\xi_1^\nu+b^\nu)
  =\lim_{\nu\to\infty}(a^\nu\xi_2^\nu+b^\nu)=0,
\]
其中 $\lim_{\nu\to\infty}\xi_1^\nu=x_1\ne x_2
=\lim_{\nu\to\infty}\xi_2^\nu$。所以 $a^\nu$ 与 $b^\nu$ 都趋于零。
这意味着 $\psi^\nu$ 在
$\mathbb C=S^2\setminus\{\infty\}$ 的紧子集上\BookRecordIndex{T0416}逐紧一致收敛于
$y_0=0$。由于 $\sigma^\nu$ 和 $\rho^\nu$ 都在整个 $S^2$ 上一致收敛于
恒等变换，$\phi^\nu$ 也在 $\mathbb C$ 的紧子集上\BookRecordIndex{T0416}逐紧一致收敛于零。
这证明了引理~\ref{lem:D-1-4}。
\end{proof}
